\section{The curvature factor and the retained horizon}
\label{sec:curvature-horizon-producer}
The positive beds of Section~\ref{sec:positive-row-producer} must
be joined by one common family of transports. The remaining issue
is the positivity of the curvature defect. The construction first
writes that defect as a double exterior derivative, then adds a
fixed positive amendment before taking its square root.

\subsection{An exponentially summable exterior lift}
We give the row-level lifting argument, because a cycle on a
single finite torus would not suffice. Write
$q_a(p)=e^{ip_a}-1$ and $d_r(p)=q(p)\wedge\cdot$ on $r$-forms.
Let $A(p)$ be the symbol of an exponentially summable row from
$(r+1)$-forms to a fixed coordinate space, where $r+1>0$, and
suppose $A(p)d_r(p)=0$ for every real $p$.

For one coordinate $a$, let $A_0(p)=A(p_1,\ldots,0,\ldots,p_4)$.
The difference $A-A_0$ is divisible by $q_a$ as a Laurent row.
To see this coefficientwise, replace each monomial in that
coordinate by
\[
 \frac{e^{inp_a}-1}{e^{ip_a}-1}
 =\begin{cases}
   \sum_{h=0}^{n-1}e^{ihp_a},& n>0,\\
   0,&n=0,\\
   -\sum_{h=n}^{-1}e^{ihp_a},&n<0.
  \end{cases}
\]
These are finite identities, including at $p_a=0$ by continuity.
If the original row has weighted norm at width $\eta$, the
quotient has weighted norm at width $\eta/2$ bounded by
$2/\eta$ times the original one. Each monomial produces $|n|$
terms of no greater absolute displacement, and
\[
       |n|e^{(\eta/2)|z|_1}
                  \le (2/\eta)e^{\eta|z|_1}
                   \quad(|n|\le|z|_1).
\]
This proves both convergence and the claimed bound.

Let $E_a$ denote exterior multiplication by the $a$th basis
vector, and $\iota_a=E_a^{\mathsf T}$ the corresponding
contraction. The exterior identity
$d_r\iota_a+\iota_a d_{r+1}=q_a I$ follows on every coordinate
face by distinguishing whether that face contains $a$; the two
possible orders of adjoining another coordinate have opposite
signs. Set
\[
       R_a=\frac{A-A_0}{q_a}\,\iota_a,
       \qquad S_a=A_0(I-E_a\iota_a).
\]
Then
\begin{equation}\label{eq:curvature-axis-lift}
                            R_a d_{r+1}+S_a=A.
\end{equation}
For $q_a\ne0$, multiply this identity by $q_a$ and use
$Ad_r=0$ and $A_0d_r=q_aA_0E_a$; the latter follows from
the cycle identity at $p_a=0$. This proves it there, and
continuity proves it also at $q_a=0$. Multiplication on the
right by $d_r$ shows that $S_a$ remains a cycle.
The row $S_a$ has zero columns on faces containing $a$, and
on the other columns it simply evaluates $p_a$ at zero.

Apply this operation successively in the four coordinate
directions. Previously vanished columns remain zero. Every
positive-degree face contains some coordinate, so the final
residual is zero. Summing the four successive $R_a$ gives a
row $\mathcal L(A)$ with
\[
                         \mathcal L(A)d_{r+1}=A.
\]
There are only four operations. The width is reduced a fixed
number of times and the cost is multiplied by a fixed constant;
any scalar scale factor in the initial bound is retained.
Transposition gives the left-handed lift. All these identities
hold on every finite period by Fourier diagonalization.
Conversely, an exponentially summable row identity valid on all
finite periods holds for all real symbols: the finite Fourier
frequencies are dense and the symbols are continuous.

We shall also need a two-sided version. Suppose a two-form row
$E$ satisfies $Ed_1=0$ and $d_1^{\mathsf T}E=0$ on every
period. Let $W$ be a left lift, so $d_2^{\mathsf T}W=E$,
and put $V=Wd_1$. Then $d_2^{\mathsf T}V=0$. Lift once more
on the left to obtain $T$ with $d_3^{\mathsf T}T=V$.
The row $Td_0$ is annihilated on the left by $d_3^{\mathsf T}$.
A further left lift would take values in five-forms, whose
coordinate space is zero in dimension four. Hence $Td_0=0$.
The right lift therefore gives $X d_1=T$. Consequently
\[
       W'=W-d_3^{\mathsf T}X
       \quad\hbox{satisfies}\quad
       W'd_1=0,\qquad d_2^{\mathsf T}W'=E.
\]
A final right lift $C$ with $Cd_2=W'$ gives
\begin{equation}\label{eq:curvature-double-lift}
                         d_2^{\mathsf T}Cd_2=E.
\end{equation}
If $E$ is symmetric, replace $C$ by
$(C+C^{\mathsf T})/2$. The identity and exponential row bounds
are preserved, and the new $C$ is symmetric.

\subsection{The unamended defect and its scale factor}
For one flight let $H_c,B_c$ and $H_f,B_f$ be the reciprocal
pulse and dual-pulse rows on the coarse and fine lattices.
Let $T$ be its two-form transport and $A_0$ its one-form
transport. The conditional mean is denoted by $Q$. The exact
unamended transport identities are
\begin{equation}\label{eq:curvature-tangent}
 T^{\mathsf T}d_{1,f}=d_{1,c}Q,\qquad
 TH_cd_{1,c}=H_fd_{1,f}A_0,\qquad
 QA_0=I-k^{-1}d_{1,c}^{\mathsf T}H_cd_{1,c}.
\end{equation}
These identities were proved for the original kernel in
Section~\ref{sec:kernel-transport-identities};
the amendment does not alter them on exact forms, by
\eqref{eq:positive-amendment-intertwine} and $Gd_1=Hd_1$.
Define
\[
 Z_0=B_c-T^{\mathsf T}B_fT-k^{-1}d_{1,c}d_{1,c}^{\mathsf T},
                        \qquad E=H_c Z_0 H_c.
\]
The matrices are symmetric. Moreover
\[
 T^{\mathsf T}B_fT H_cd_{1,c}
       =T^{\mathsf T}d_{1,f}A_0
       =d_{1,c}\bigl(I-k^{-1}d_{1,c}^{\mathsf T}H_cd_{1,c}\bigr).
\]
Substituting this in $Z_0H_cd_{1,c}$ gives zero. Thus
$Ed_{1,c}=0$, and symmetry gives the left cycle identity.
The two-sided lift yields a symmetric row $C$ such that
\begin{equation}\label{eq:curvature-unamended-factor}
 Z_0=(d_{2,c}B_c)^{\mathsf T}C(d_{2,c}B_c).
\end{equation}
Indeed multiply \eqref{eq:curvature-double-lift} by $B_c$ on
both sides and use $B_cH_c=H_cB_c=I$.

We check uniformity in the block side $L$. The dual transport
and the row $B_fT$ have entrywise strip cost $O(L^{-2})$,
including their required cover differences, by the wire estimates
of Sections~\ref{sec:ex-lattice-free-kernel-estimates} and
\ref{sec:ex-lattice-free-kernel-components}. The Gram product
$T^{\mathsf T}B_fT$ sums over $L^4$ fine residues and a fixed
number of form labels. Its entrywise cost is consequently
$O(L^4 L^{-2}L^{-2})=O(1)$. The other rows in $E$ have uniform
costs. All lifts use fixed exterior index sets on the coarse
side, so the resulting cost of $C$ is uniform in $L$ too.

There is a second lift with a stronger scale gain. Put
\[
                         Y=d_{2,f}B_fTH_c.
\]
The second identity in \eqref{eq:curvature-tangent} gives
$Yd_{1,c}=d_{2,f}d_{1,f}A_0=0$. Its right lift is a row $S$
with $Sd_{2,c}=Y$. One fine exterior difference is a signed
sum of at most four cover differences of $B_fT$; the first
difference estimate supplies one additional $L^{-1}$.
Thus each entry of $Y$, and then of $S$, has strip cost
$O(L^{-3})$. Multiplying the lift identity by $B_c$ gives
\begin{equation}\label{eq:curvature-charge-lift}
                     S d_{2,c}B_c=d_{2,f}B_fT.
\end{equation}
The row $S^{\mathsf T}S$ now has cost $O(L^{-2})$:
its summation has $L^4$ residues, while both factors cost
$L^{-3}$. Retaining this count is essential.

\subsection{Making the curvature row positive}
Choose amendment coefficients $m,l\ge0$ for the fine and
coarse beds. With $D_c=(I+l d_{2,c}B_cd_{2,c}^{\mathsf T})^{-1}$,
set
\[
                R_0=C+lI-mS^{\mathsf T}S,
                       \qquad R=D_c^{\mathsf T}R_0D_c.
\]
The uniform bounds just proved give constants $C_0,C_1>0$,
independent of the train and $L$, with
\[
 |\langle u,Cu\rangle|\le C_0\|u\|^2,
 \qquad
 \langle u,S^{\mathsf T}Su\rangle
                                 \le C_1L^{-2}\|u\|^2.
\]
Choose $l\ge C_0+2$ and $L^2\ge mC_1$. Then
\[
                         \langle u,R_0u\rangle\ge\|u\|^2.
\]
This is positivity on the entire three-form space. No projection
onto a momentum-dependent range is needed. The row
$D_c^{-1}=I+l d_{2,c}B_cd_{2,c}^{\mathsf T}$ has uniformly
bounded operator norm, say at most $K$. Therefore
\[
           \langle u,Ru\rangle\ge\|D_cu\|^2
                                      \ge K^{-2}\|u\|^2.
\]
The products defining $R$ have a uniform exponential row bound.
Section~\ref{sec:positive-row-producer} supplies its positive
square-root row $P$, uniformly over all flights of these stages.

It remains to identify this factor with the actual amended
curvature defect. Write
\[
 A_c=B_c+l(d_{2,c}B_c)^{\mathsf T}(d_{2,c}B_c),\qquad
 A_f=B_f+m(d_{2,f}B_f)^{\mathsf T}(d_{2,f}B_f),
 \qquad G_c=A_c^{-1}.
\]
Equations~\eqref{eq:curvature-unamended-factor} and
\eqref{eq:curvature-charge-lift}, followed by expansion, give
\[
 A_c-T^{\mathsf T}A_fT-k^{-1}d_{1,c}d_{1,c}^{\mathsf T}
           =(d_{2,c}B_c)^{\mathsf T}R_0(d_{2,c}B_c).
\]
Using $(d_{2,c}B_c)G_c=D_cd_{2,c}$ from
\eqref{eq:positive-amendment-intertwine}, multiply by $G_c$
on both sides to obtain
\begin{equation}\label{eq:curvature-paid-factor}
 G_c\bigl(A_c-T^{\mathsf T}A_fT
               -k^{-1}d_{1,c}d_{1,c}^{\mathsf T}\bigr)G_c
                         =d_{2,c}^{\mathsf T}P^{\mathsf T}P d_{2,c}.
\end{equation}
This is precisely the curvature factorization in
Section~\ref{sec:ex-lattice-rg-horizon-geometry}, with the actual
amended beds. All formulas hold on every finite period, since
they were constructed from the same Laurent rows.

\subsection{Choosing the prices before the scale sequence}
For each adjacent pair of stages $j,t$, where $t=j$ or $t=j-1$,
let $a_{j,t}>0$ be the lower requirement on the coarse amendment
coefficient just obtained. Select
\[
                         l_i=\max(a_{i,i},a_{i+1,i}).
\]
Then every permitted transition has $l_t\ge a_{j,t}$ and
$l_j\ge0$. Its required block-size floor is finite because the
fine amendment $m=l_j$ is now fixed. Combine the two permitted
flight floors at stage $j$ by their maximum. Combine bed and
flight strip widths and positive lower bounds by minima, and
upper bounds by maxima. All widths and lower bounds stay positive,
since each combination is finite.

Choose the stage's direct precision below its positive bed
ceiling. Choose its block length above the flight floor. These
are requirements on the original stage profiles and pre-length
floors; they are imposed before the common schedule of
Section~\ref{sec:common-scale-selection}. They hold for both
within-stage steps and a change to the next stage, however many
times either occurs.

At a terminal cutoff $N$, start with the top train and descend
along those very block lengths. At each level use its positive
bed from Section~\ref{sec:positive-row-producer}; at each step
use \eqref{eq:curvature-paid-factor} and the exact block completion
of Section~\ref{sec:ex-lattice-rg-horizon-geometry}. The latter's
strip calculation supplies the prescribed iterated difference bounds
for its unchanged $A_0,T$ rows and the displayed strip bounds for
its remaining blocks. This changes neither the bed rows nor the
original one-form and two-form transports $A_0,T$. It only
constructs the remaining three blocks from their common positive
curvature factor. Taking the indicated finite minima and maxima at each
stage gives stage-indexed prices valid for every $N$, every intermediate
level and every finite period. These prices are fixed before the level
schedule, but need not have a common positive lower width or a common
upper cost over all stages. This proves the bed
and curvature suppliers used by the horizon construction, using
the original pulse transport identities proved in
\eqref{eq:transport-installed-identities}. The attachment and bounds of the
nonlinear path remain the separate construction step.
